\section{Lubin--Tate formal group laws}

Let $K$ be a field complete with respect to a discrete valuation, with finite
residue field. For instance, $K$ is a nonarchimedean local field. Let $\calO_K$ be
the unit ball of $K$, let $q$ be the cardinality of the residue
field, and let $\pi$ be a \emph{uniformizer} of $\calO_K$.
A uniformizer of $K$ is an element $\pi \in \mathcal O_K$ such that
$\gothm_K = (\pi)$.
Equivalently, if $v_K$ is the normalized discrete valuation on $K$, then
$\pi$ is a uniformizer if $v_K(\pi)=1$.
Thus every nonzero element $x\in K^\times$ can be written uniquely as
$x = u\pi^n$, where $u\in \mathcal O_K^\times$ is a unit and $n\in \mathbb Z$.

\begin{definition}[{\cite[Chap.~6 Definition 3.3]{cassels1967algebraic}}]
    Let $\calF_\pi$ denote the set of power series $f(X) \in \calO_K[\![X]\!]$ which
    satisfy the two conditions:
    \begin{align}
        f(X) &\equiv \pi X \pmod{\deg 2}, \\
        f(X) &\equiv X ^ q \pmod{\pi}.
        \label{Lubin_Tate_condition}
    \end{align}

\end{definition}

\begin{lemma}[{\cite[Chap.~6 Proposition 5]{cassels1967algebraic}}]
    Let $f(X)$ and $g(X)$ be elements of $\calF_\pi$ and let $L(X_1, \ldots, X_n)
    = \sum_{i = 1}^n a_i X_i$ be a linear form with coefficients in $\calO_K$.

    Then there
    is a unique power series $F(X_1, \ldots, X_n)$ over $\calO_K$ such that
    \begin{align}
        F (X_1, \ldots, X_n) &\equiv  L(X_1, \ldots, X_n) \pmod{\deg 2}, \label{eq: LT_trunc_two} \\
        f(F(X_1, \ldots, X_n)) &= F(g(X_1), \ldots, g(X_n)). \label{eq: LT_subst}
    \end{align}
    \label{lemma: constructive_lemma_lubin_Tate}
\end{lemma}

\begin{proof}
    % First, we claim that for all $k \in \NN$,
    % there is a polynomial $F_k(X_1, \ldots, X_n)$ of degree $k$ such that

    % \begin{align*}
    %     F_k(X_1, \ldots , X_n) &\equiv L (X_1, \ldots, X_n) \pmod{\deg 2} \quad \text{and} \\
    %     f(F_k(X_1, \ldots , X_n)) &\equiv F_k(g(X_1), \ldots, g(X_n)) \pmod{\deg (k + 1)}.
    % \end{align*}

    % Then we obtain a sequence of polynomial $F_k(X_1, \ldots, X_n)$. Define
    % $F$ to be $\lim F_k$. Then the $F$ satisfies the equation \eqref{eq: LT_trunc_two}
    % and \eqref{eq: LT_subst}. Indeed, for equation \eqref{eq: LT_trunc_two}, it follows
    % that $F_k(X_1, \ldots , X_n) \equiv L (X_1, \ldots, X_n) \pmod{\deg 2}$ for all $k$. For
    % equation \eqref{eq: LT_subst}, we have that for all $k$,
    % \[f(F(X_1, \ldots, X_n)) \equiv f(F_k(X_1, \ldots, X_n)) \equiv
    % F_k(g(X_1), \ldots, g(X_n)) \pmod{\deg (k + 1)}. \]

    % The first congruence follows from $f$ is of constant coefficient zero, and
    % \begin{equation}
    %     F(X_1, \ldots, X_n) \equiv F_k(X_1, \ldots, X_n) \pmod{\deg (k+1)}.
    %     \label{eq: LT_trunc_k}
    % \end{equation}

    % Moreover, we have that

    % \[F(g(X_1), \ldots, g(X_n)) \equiv F_k(g(X_1), \ldots, g(X_n)) \pmod{\deg (k + 1)}. \]

    % Then we conclude that for all $k$,
    % \[f(F(X_1, \ldots, X_n)) \equiv F(g(X_1), \ldots, g(X_n)) \pmod{\deg (k+1)}. \]

    % It follows that $f(F(X_1, \ldots, X_n)) = F(g(X_1), \ldots, g(X_n))$.

    % \textit{Proof of the claim:} For $k = 1$, we take $F_1(X_1, \ldots, X_n) = L(X_1, \ldots, X_n)$.

    % Now assume that the result holds for $k$.
    % Write
    % \[f (F_k(X_1,\ldots, X_n )) \equiv
    %     F_k(g(X_1), \ldots, g(X_n)) + E_{k+1} \pmod{\deg (k+ 2)}.\]

    % According to the induction hypothesis, $E_{k+1}$ is homogeneous of
    % degree $k + 1$. Define \[F_{k+1}(X_1, \ldots, X_n) \coloneqq F_k(X_1, \ldots, X_n) -
    % E_{k+1}/ (\pi - \pi ^{k+1}).\]

    % % On the other hand, since the uniqueness of $F_k(X_1, \ldots, X_n)$, then
    % % $F_{k+1}(X_1,\ldots,X_n)$ should be of the form $F_{k+1}(X_1,\ldots, X_n) =
    % %     F_k(X_1, \ldots, X_n) + \phi_{k+1}(X_1, \ldots, X_n)$, where
    % %     $\phi_{k+1}(X_1, \ldots, X_n)$ is homogeneous of degree $(k+1)$, due to the definition of
    % %     $\calF_\pi$.

    % Since $f(X) \equiv \pi X \pmod{\deg 2}$ and $g(X) \equiv \pi X \pmod{\deg 2}$, then we have
    % \begin{align*}
    %     f(F_{k+1}(X_1, \ldots, X_n)) &\equiv F_k(g(X_1), \ldots, g(X_n)) -
    %     \pi E_{k+1}/ (\pi - \pi ^{k+1}) \pmod{\deg (k + 2)}, \\
    %     F_{k+1}(g(X_1), \ldots, g(X_n)) &\equiv  F_k(g(X_1), \ldots, g(X_n)) -
    %     \pi ^ {k+1} E_{k+1}/ (\pi - \pi ^{k+1})\pmod{\deg (k + 2)}.
    % \end{align*}

    % The last congruence follows from the fact that for any $\phi$ is homogeneous of degree $(k+1)$,
    % that is,
    % \[\phi (g(X_1), \ldots, g(X_n)) \equiv \phi(\pi X_1, \ldots, \pi X_n)
    % \equiv \pi ^ {k+1} \phi (X_1, \ldots, X_n) \pmod{\deg (k + 1)}. \]

    % Therefore, we conclude that
    % \[E_{k+1} + (\pi- \pi^{k+1}) \phi_{k+1} (X_1, \ldots, X_n) = 0.\]

    % Then $\phi_{k+1}(X_1, \ldots, X_n) = - E_{k+1}/ (\pi - \pi^{k+1})$.

    First, we claim that for all $k \in \NN$,
    there is a unique polynomial $F_k(X_1, \ldots, X_n)$ of degree $k$ such that

    \begin{align*}
        F_k(X_1, \ldots , X_n) &\equiv L (X_1, \ldots, X_n) \pmod{\deg 2} \quad \text{and} \\
        f(F_k(X_1, \ldots , X_n)) &\equiv F_k(g(X_1), \ldots, g(X_n)) \pmod{\deg (k + 1)}.
    \end{align*}

    Then we obtain a sequence of polynomials $F_k(X_1, \ldots, X_n)$. Define
    $F$ to be $\lim F_k$. Then $F$ satisfies the equation~\eqref{eq: LT_trunc_two}
    and~\eqref{eq: LT_subst}. Indeed, for equation~\eqref{eq: LT_trunc_two}, it follows
    that $F_k(X_1, \ldots , X_n) \equiv L (X_1, \ldots, X_n) \pmod{\deg 2}$ for all $k$. For
    equation~\eqref{eq: LT_subst}, we have that for all $k$,
    \[f(F(X_1, \ldots, X_n)) \equiv f(F_k(X_1, \ldots, X_n)) \equiv
    F_k(g(X_1), \ldots, g(X_n)) \pmod{\deg (k + 1)}. \]

    The first congruence follows from the fact that $f$ has constant coefficient zero, together with
    \begin{equation}
        F(X_1, \ldots, X_n) \equiv F_k(X_1, \ldots, X_n) \pmod{\deg (k+1)}.
        \label{eq: LT_trunc_k}
    \end{equation}

    Moreover, we have that

    \[F(g(X_1), \ldots, g(X_n)) \equiv F_k(g(X_1), \ldots, g(X_n)) \pmod{\deg (k + 1)}. \]

    Then we conclude that for all $k$,
    \[f(F(X_1, \ldots, X_n)) \equiv F(g(X_1), \ldots, g(X_n)) \pmod{\deg (k+1)}. \]

    It follows that $f(F(X_1, \ldots, X_n)) = F(g(X_1), \ldots, g(X_n))$.
    Finally, the uniqueness in this proposition follows from equation~\eqref{eq: LT_trunc_k} and the uniqueness in
    the claim.

    \textit{Proof of the claim:} For $k = 1$, we take $F_1(X_1, \ldots, X_n) = L(X_1, \ldots, X_n)$. \par
    Now assume that the result holds for $k$. Write
    \[f (F_k(X_1,\ldots, X_n )) \equiv
        F_k(g(X_1), \ldots, g(X_n)) + E_{k+1} \pmod{\deg (k+ 2)}.\]

    According to the induction hypothesis, $E_{k+1}$ is homogeneous of
    degree $k + 1$.

    On the other hand, by the uniqueness of $F_k(X_1, \ldots, X_n)$,
    $F_{k+1}(X_1,\ldots,X_n)$ should be of the form $F_{k+1}(X_1,\ldots, X_n) =
        F_k(X_1, \ldots, X_n) + \phi_{k+1}(X_1, \ldots, X_n)$, where
        $\phi_{k+1}(X_1, \ldots, X_n)$ is homogeneous of degree $(k+1)$, due to the definition of
        $\calF_\pi$. Then we have
    \begin{align*}
        f(F_{k+1}(X_1, \ldots, X_n)) &\equiv F_k(g(X_1), \ldots, g(X_n)) +
        \pi \phi_{k+1}(X_1, \ldots, X_n) \pmod{\deg (k + 2)}, \\
        F_{k+1}(g(X_1), \ldots, g(X_n)) &\equiv  F_k(g(X_1), \ldots, g(X_n)) +
        \pi ^ {k+1} \phi_{k+1} (X_1, \ldots, X_n) \pmod{\deg (k + 2)}.
    \end{align*}

    The last congruence follows from the fact that $\phi_{k+1}$ is homogeneous of degree $(k+1)$,
    that is,
    \[\phi_{k+1} (g(X_1), \ldots, g(X_n)) \equiv \phi_{k+1}(\pi X_1, \ldots, \pi X_n)
    \equiv \pi ^ {k+1} \phi_{k+1} (X_1, \ldots, X_n) \pmod{\deg (k + 1)}. \]

    Therefore, we conclude that
    \[E_{k+1} + (\pi- \pi^{k+1}) \phi_{k+1} (X_1, \ldots, X_n) = 0.\]

    Note that $1 - \pi ^k$ is invertible in $\calO_K$. And
    \[E_{k+1} \equiv f(F_k(X_1, \ldots, X_n)) - F_k(g(X_1), \ldots, g(X_n))
    \equiv F_k^q - F_k(X_1 ^ q, \ldots, X_n^q) \equiv 0
    \pmod{\pi}.  \]
    Then $\phi_{k+1}(X_1, \ldots, X_n) = - E_{k+1}/ (\pi - \pi^{k+1})$.
\end{proof}

\begin{definition}
    For any $f \in \calF_\pi$, we let $F_f(X,Y)$ be the unique solution of
    \begin{align*}
        F_f(X,Y) &\equiv X + Y \pmod{\deg 2}, \\
        f(F_f(X,Y)) &= F_f(f(X), f(Y)).
    \end{align*}
\end{definition}


\begin{theorem}[Lubin--Tate Formal Group Laws {\cite[Chap.~VI.3.3 Proposition 1]{cassels1967algebraic}}]
    For power series $f \in \calF_\pi$, we have
    \begin{align}
        F_f(X,Y) &= F_f(Y,X) \label{eq: lubin_Tate_comm} \\
        F_f(F_f(X,Y),Z) &= F_f(X, F_f(Y,Z)) \label{eq: lubin-Tate_assoc}
    \end{align}
     As a result, $F_f(X,Y)$ is a commutative formal group law over $\calO_K$.

    \label{thm: Lubin_Tate}
\end{theorem}

\begin{proof}
    For equation~\eqref{eq: lubin_Tate_comm}, let $G (X,Y)= F(Y,X)$. Then $G(X,Y)
    \equiv X+Y \pmod{\deg 2}$ and
    \[f(G(X,Y)) = f(F(Y,X)) = F(f(Y),f(X)) = G(f(X),f(Y)). \]
    Therefore, by the uniqueness in Lemma~\ref{lemma: constructive_lemma_lubin_Tate},
    $F(X,Y) = G(X,Y) = F(Y,X)$.

    For equation~\eqref{eq: lubin-Tate_assoc}, let
    \[ G_1(X,Y,Z) = F_f(F_f(X,Y),Z) \quad \text{and} \quad
    G_2(X,Y,Z) = F_f(X,F_f(Y,Z)).\]
    We can check that $G_i(X,Y,Z) \equiv X+Y+Z \pmod{\deg 2}$ and
    \begin{equation*}
        f(G_i(X,Y,Z)) = G_i(f(X),f(Y),f(Z)).
    \end{equation*}
    Indeed,
    \[f(G_1(X,Y,Z)) = f(F_f(F_f(X,Y), Z)) = F_f(f(F_f(X,Y)), f(Z)) =
    F_f(F_f(f(X), f(Y)), f(Z)). \]
    \[f(G_2(X,Y,Z)) = f(F_f(X, F_f(Y,Z))) = F_f(f(X), f(F_f(Y,Z))) =
    F_f(f(X), F_f(f(Y), f(Z))).  \]
    Then, by uniqueness in Lemma~\ref{lemma: constructive_lemma_lubin_Tate},
    we conclude the desired equation~\eqref{eq: lubin-Tate_assoc}.
\end{proof}


\begin{definition}
    For each $a \in \calO_K$ and $f, g \in \calF_\pi$, let $[a]_{f,g} (X)$ be the unique
    solution of
    \begin{align}
        [a]_{f,g} (X) &\equiv aX \pmod{\deg 2}, \\
        f([a]_{f,g}(X)) &= [a]_{f, g}(g(X)). \label{eq: LT_smul}
    \end{align}
    In particular, we write $[a]_{f}$ instead of $[a]_{f,f}$ to simplify notation.
\end{definition}

\begin{remark}
    Note that for any $f\in \calF_\pi$,
    \[f(X) \equiv \pi X \equiv [\pi]_f \pmod{\deg 2} \quad  \text{and} \quad
    f \circ [\pi]_f = [\pi]_f \circ f. \]
    Then by uniqueness in Lemma~\ref{lemma: constructive_lemma_lubin_Tate},
    it follows that $[\pi]_f = f$.
\end{remark}

\begin{theorem}[{\cite[Chap.~VI 3.3 Proposition 2]{cassels1967algebraic}}]
    For each $a \in \calO_K$ and $f, g \in \calF_\pi$, the power series $[a]_{f,g}(X)$ is a
    formal group homomorphism from $F_g(X,Y)$ to $F_f(X,Y)$, i.e.
    \[[a]_{f,g} F_g(X,Y) = F_f([a]_{f,g}(X), [a]_{f,g} (Y)).\]
\end{theorem}

\begin{proof}
    For simplicity, we denote $[a]_{f,g}$ by $h$. We have $h(F_g(X,Y)) \equiv aX+aY
    \pmod{\deg 2}$, $F_f(h(X), h(Y)) \equiv aX+ aY \pmod{\deg 2}$ and
    \begin{align*}
        h\circ F_g(g(X),g(Y)) &= h\circ g(F_g(X,Y)) = f \circ h (F_g(X,Y)) \\
        F_f(h(g(X)), h(g(Y))) &= F_f(f(h(X)), f(h(Y))) = f \circ F_f(h(X), h(Y))
    \end{align*}
    Then the result follows from uniqueness in Lemma~\ref{lemma: constructive_lemma_lubin_Tate}.
    \end{proof}

\begin{theorem}[{\cite[Theorem 1]{LubinTate1965}}]
    Let $a, b \in \calO_K$, $f, g, h \in \calF_\pi$, then
    \[[a]_{f,g} ([b]_{g,h}(X)) = [ab]_{f,h} (X).\]
\end{theorem}

\begin{proof}
    We have $[a]_{f,g} ([b]_{g,h}(X)) \equiv abX \pmod{\deg 2}$ and
    \begin{equation*}
        f \circ ([a]_{f, g} \circ [b]_{g,h}) = [a]_{f,g} \circ (g \circ [b]_{g,h})
        = [a]_{f,g}\circ [b]_{g,h} \circ h.
    \end{equation*}
    Thus, $[ab]_{f,h}$ and $[a]_{f,g} \circ [b]_{g,h}$ are solutions satisfying these conditions.
    By uniqueness in Lemma~\ref{lemma: constructive_lemma_lubin_Tate}, they coincide.
\end{proof}

\begin{theorem}[{\cite[Theorem 1]{LubinTate1965}}]
    Let $a, b \in \calO_K$, $f, g \in \calF_\pi$, then
    \[F_f([a]_{f,g}(X), [b]_{f,g}(X)) = [a+b]_{f,g} (X).\]
\end{theorem}

\begin{proof}
    Note that $F_f([a]_{f,g}(X), [b]_{f,g}(X)) \equiv (a+b)X \pmod{\deg 2}$ and
    \begin{equation*}
        f(F_f([a]_{f,g}(X), [b]_{f,g}(X))) = F_f(f\circ [a]_{f,g}(X), f\circ [b]_{f,g}(X))
            = F_f([a]_{f,g}\circ f(X), [b]_{f,g}\circ f(X)).
    \end{equation*}
    Therefore, by the uniqueness in Lemma~\ref{lemma: constructive_lemma_lubin_Tate}, we have
    \[ F_f([a]_{f,g}(X), [b]_{f,g}(X)) = [a+b]_{f,g} (X).\]
\end{proof}

\begin{theorem}[{\cite[Chap.~VI 3.3 Proposition 4]{cassels1967algebraic}}]
    Let $f,g\in \calF_\pi$, then $F_f$ is isomorphic to $F_g$, and there exists a unique strong
    isomorphism from $F_f$ to $F_g$. This means that the Lubin--Tate formal group $F_f$ is
    independent of $f \in \calF_\pi$, and depends only on $\pi$.
    \label{thm: Lubin_Tate_independet_f_g}
\end{theorem}

\begin{proof}
    Take a unit $a \in \calO_K$, then
    \[[a]_{f,g} \circ [a^{-1}]_{g,f} = [a a ^ {-1}]_f = [1]_f. \quad \text{and} \quad
    [a ^ {-1}]_{g,f} \circ [a]_{f,g} = [a^{-1}a]_g = [1]_g.
    \]
    It follows that $[a]_{f,g}$ is a formal group isomorphism from $F_g$ to $F_f$ with
    its inverse $[a^{-1}]_{g,f}$.
    And the unique strong isomorphism is
    $[1]_{f,g}$.
\end{proof}

\begin{remark}
    Let $L$ be an algebraic extension of $K$, let $\gothm_L$ be the maximal ideal in the
    ring of integers of $L$. Given elements $x_1, \ldots, x_n \in \gothm_L$ and a formal
    series $G(X_1, \ldots, X_n) \in \calO_K[\![X_1, \ldots, X_n]\!]$, the series
    $G(x_1, \ldots, x_n)$ converges to an element in $\gothm_L$ if the constant term of
    $G(X_1, \ldots, X_n)$ is zero. By the theorem above, for each power series $f\in \calF_\pi$, this
    determines an $\calO_K$-module structure on the set $\gothm_L$. More precisely, addition
    and scalar multiplication are defined by
    \begin{align*}
        x + y &= F_f(x,y),  \\
        ax &= [a]_f (x).
    \end{align*}
    \label{rem: A-module}
\end{remark}

\section{Applications}
Fix an algebraic closure $\overline{K}$ of $K$, and let $\Fsep{K}$ be a separable closure of $K$.
Given $f\in \calF_\pi$, for each integer $n \ge 1$, let $f^n(X)$ be the $n$-th iterate of $f(X)$.
Let $\Lambda_{f,n}$ be the set of all roots of $f^n(X)$ in $\overline{K}$ and let
$K(\Lambda_{f,n})$ be the field generated over $K$ by these elements.
Denote $\Lambda_f = \bigcup_{n=1}^\infty \Lambda _{f,n} $.

\begin{lemma}
    For any $f,g\in \calF_\pi$, $\lambda \in \Lambda_{f,n}$ if and only
    if $[1]_{g,f} (\lambda) \in \Lambda_{g,n}$.
\end{lemma}

\begin{proof}
    By equation~\eqref{eq: LT_smul} and induction, we have that for
    any $n \in \NN_{>0}$,
    \[g^n \circ [1]_{g,f} = [1]_{g,f} \circ f^n. \]

    Suppose that $\lambda \in \Lambda_{f, n}$, so that $f^n(\lambda) = 0$.
    Then $g^n([1]_{g,f}(\lambda)) = [1]_{g,f} (f^n(\lambda)) = [1]_{g,f}(0) = 0$.
    We conclude that $[1]_{g,f} (\lambda) \in \Lambda_{g,n}$.

    Suppose $[1]_{g,f}(\lambda) \in \Lambda_{g, n}$. Denote
    $[1]_{g,f}(\lambda)$ by $\mu$. Then by the discussion above, we have that
    $[1]_{f,g}(\mu) \in \Lambda_{f, n}$. Note that
    \[[1]_{f,g}(\mu) = [1]_{f,g}([1]_{g,f}(\lambda)) = \lambda. \]

    Then we conclude that $\lambda \in \Lambda_{f, n}$.
\end{proof}

\begin{remark}

    It follows that the field $K(\Lambda_{f,n})$ is independent of $f$ and depends only
    on the uniformizer $\pi$. Therefore, we denote $K(\Lambda_{f,n})$ by
    $K_{\pi,n}$. Note that
    \[K(\Lambda_f) = K \left(\bigcup_{n=1}^{\infty} \Lambda_{f, n}\right) =
    \bigcup_{n=1}^{\infty} K (\Lambda_{f,n}).\]

    The second equality follows from the fact that
    $\{ K(\Lambda_{f, n}) \}_n$ is an increasing chain of fields.
    It follows that $K(\Lambda_f)$ is independent of $f$, so we denote $K(\Lambda_f)$ by $K_\pi$.
    Moreover, we denote $\Gal(K_{\pi,n}/K)$ by $G_{\pi, n}$
    and $\varprojlim G_{\pi,n}$ by $G_\pi$.
\end{remark}

According to Remark~\ref{rem: A-module}, $F_f$ defines an $\calO_K$-module structure on
$\gothm_{\Fsep{K}}$ and we denote this $\calO_K$-module as $M_f(\Fsep{K})$.
Then $\Lambda_{f,n}$ consists of those
elements $\lambda \in \gothm_{\Fsep{K}}$ such that $\pi ^n \lambda = 0$.

\begin{theorem}[{\cite[Chap.~VI 3.3 Theorem 3]{cassels1967algebraic}}]
    $\Lambda_{f,n}$ is isomorphic to $\calO_K/ \pi^n \calO_K$ as an $\calO_K$-module and
    $\Lambda_f$ is isomorphic to $K / \calO_K $ as an $\calO_K$-module.
    \label{thm: Lam_iso_K_quot_O}
\end{theorem}

In order to prove this theorem, we need a result about divisible modules over a PID.

\begin{definition}[{\cite{anderson1992rings}, Chap.~5 Definition 18.10}]
Let $R$ be a ring and $S$ be a simple left $R$-module. An \emph{injective hull}
(or \emph{injective envelope}) of $S$, denoted by $E(S)$, is an injective $R$-module $E$
together with an \emph{essential monomorphism} $\iota \colon S \to E$. Here, \emph{essential} means
that every nonzero submodule of $E(S)$ has nonzero intersection with $S$.

Equivalently, $E(S)$ can be characterized as:
\begin{enumerate}
    \item The maximal essential extension of $S$.
    \item The minimal injective extension of $S$.
\end{enumerate}
\end{definition}

\begin{remark}
    The injective hull $E(S)$ of a simple module $S$ is necessarily an
    \emph{indecomposable} injective module; that is, $E(S)$ cannot be written as
    a direct sum of two nonzero submodules.
\end{remark}

\begin{theorem}[Structure theorem of divisible modules over a PID, {\cite{matlis1958injective},
    Chap~10, Prop~3.1}]
    Let $R$ be a principal ideal domain and $K$ be its field of fractions. Then every divisible
    $R$-module is isomorphic to a direct sum of $K$ and copies of the injective hulls
    of the simple module $R/ (p)$, varying over prime elements $p\in R$. \par
    Let $M$ be a divisible $R$-module, then $M$ can be written as follows:
    \[M \cong \left(\bigoplus_{i \in I} K\right) \oplus \left(\bigoplus_{p \ prime} \bigoplus_{j \in J_p}
        E(R / (p))\right), \quad \text{where}  \]
    \begin{itemize}
        \item $I$ is an index set measuring how many copies of the torsion-free part of $M$ are needed.
        \item $p$ varies over all prime elements in $R$.
        \item $J_p$ is an index set measuring how many copies of $p$-primary torsion part
                $E(R/(p))$ are needed for each $p$.
        \item $E(R/(p))$ is the injective hull of $R/(p)$. Over a PID, this is isomorphic to
                $K/R_{(p)}$.
    \end{itemize}
    \label{thm: structure_div_mod_PID}
\end{theorem}

% \begin{proof}
%    The proof of the structure theorem for divisible modules can be found in Chapter 3 of
%    $\cite{lam1999lectures}$ or Theorem 2.5 and Proposition 3.1 in \cite{matlis1958injective}.
% \end{proof}

\begin{proof}[Proof of theorem $\ref{thm: Lam_iso_K_quot_O}$]
    The first claim follows from the fact that $\Lambda_{f,n}$ is the $\pi^n$-torsion of
        the $\calO_K$-module $M_f(\Fsep{K})$.

    According to Theorem~\ref{thm: Lubin_Tate_independet_f_g}, we are free to choose
    $f\in \calF_\pi$ as we please. We take $f(X)= \pi X + X ^ q$. Since
    $f'(X) = q X^{q-1} + \pi \neq 0$, $\Lambda_{f, 1}$ consists of $q$ distinct roots
    of the equation $X^q + \pi X = 0 $. For any given $\alpha \in \gothm_{\Fsep{K}}$, since
    $\pi X + X ^ q = \alpha$ is separable, $\Lambda_{f,1} \subset \gothm_{\Fsep{K}}$. This means
    that for any given $\alpha \in \gothm_\Fsep{K}$, there is a $\beta$ such that
    $\pi \cdot \beta = \alpha$ in $M_f(\Fsep{K})$. It follows that for any given $r \in \calO_K$,
    $r M_f(\Fsep{K}) = M_f (\Fsep{K})$, which means that $M_f(\Fsep{K})$ is a divisible $\calO_K$-module.
    Then $\Lambda_f$ is a divisible torsion $\calO_K$-module.

    By Theorem~\ref{thm: structure_div_mod_PID} and the fact that $\calO_K$ is a PID,
    we have that
    \[\Lambda_f \cong \left(\bigoplus_{i\in I} K\right) \oplus \left(\bigoplus_{p \ prime} \bigoplus_{j \in J_p}
        E(R / (p))\right).  \]

    Since $\Lambda_f$ is a torsion module, the index $I$ should be empty, i.e.
    there are no copies of $K$ in this decomposition.
    Note that $E(R/(p))$ is isomorphic to $K/R_{(p)}$, since $\calO_K$ is a PID.
    Since $\pi$ is the only prime element in $\calO_K$ and
    $(\calO_K) _ {(\pi)} \cong \calO_K$, it is enough to determine the index set
    $J$.

    Moreover, $\Lambda_{f,1}$ is a one-dimensional vector space over the residue field $k$ (since $\Lambda_{f,1}$
    has the same cardinality as the residue field $k$).
    It follows that $|J| = \dim_k M_f(\Fsep{K})[\pi] = \dim_k \Lambda_{f,1} = 1$,
    where $J$ is the index set of copies of $K/ \calO_K$. We conclude that
    $\Lambda_f \cong K / \calO_K$.
\end{proof}

There is a natural homomorphism $\Gal (\Fsep{K}/ K) \to \End(\Lambda_f)$; this follows from
the action of the Galois group on the set of roots of a polynomial. Moreover,
since $\Lambda _f \cong K / \calO_K$ and $\End(K/\calO_K) \cong \calO_K$,
then this induces a
homomorphism $\Gal(\Fsep{K}/ K) \to \Aut(K / \calO_K) \cong U_K$.

\begin{theorem}[{\cite[Chap.~VI 3.3 Theorem 3]{cassels1967algebraic}}]
    The natural homomorphism $\Gal(\Fsep{K}/K) \to \Aut (\Lambda_f)$ is an isomorphism.
\end{theorem}

\begin{proof}
    This map is injective by definition; it remains to prove that it is surjective.
    We have an injection $\Gal(K_{\pi,n}/ K) \to U_K/ U_K^{(n)}$, where $U_K^{(n)} = 1 + \pi ^ n \calO_K$.
    Let $\alpha \in \Lambda_{f,n}$ be a primitive element, that is, an element such that
    $[\pi ^ {n - 1}]_f \alpha \neq 0$ or $\pi ^ {n - 1} \cdot \alpha \neq 0$ in $M_f(\Fsep{K})$.
    Then $\alpha$ is one of the roots of the polynomial $f^{(n)}(X)/ f^{(n-1)}(X)$,
    where $f^{(n)}(X)$ is $n$-th iterate of $f(X)$.
    Since $f(X)/X = X ^ {q -1} + \pi$, then

    \[ \frac{f^{(n)}(X)}{f ^ {(n -1)}(X)} = f^ {(n -1)}(X)^{q-1} + \pi. \]

    By the Eisenstein criterion, this is an irreducible polynomial with degree $q ^ n - q ^{n -1}$.
    Thus, according to Galois theory, the order of $\Gal(K_{\pi,n}/ K)$ is at least
    $(q-1) q ^ {n-1}$. On the other hand, $U_K/ U_K^{(n)}$ has exactly $(q-1)q ^ {n -1}$
    elements. Therefore, $\Gal(K_{\pi,n}/K) \cong U_K/ U_K^{(n)}$.
    Taking an inverse limit, it follows that
    \[ \Gal(\Fsep{K}/K) \cong \varprojlim \Gal(K_{\pi,n}/ K) \cong \varprojlim U_K / U_K^{(n)} \cong U_K. \]
\end{proof}

\begin{corollary}[{\cite[Chap.~VI 3.3 Theorem 3]{cassels1967algebraic}}]
    The element $\pi$ is a norm from $K_{\pi,n}$.

    \label{cor: norm_Lpi}
\end{corollary}

\begin{proof}
    Let $\alpha \in \Lambda_{f,n}$ be a primitive element, then $f^{(n)}(X)/ f^{(n-1)}(X)$ is a
    monic polynomial with constant coefficient $\pi$, and the norm of $-\alpha$ is $\pi$.
\end{proof}

Let $\Kur$ be the maximal unramified extension of $K$. Since $K_\pi$ is a totally ramified extension
of $K$, $K_\pi$ and $\Kur$ are linearly disjoint. As a result, the Galois group
$\Gal(K_\pi \Kur/K)$ is the direct product of $\Gal(K_\pi/K)$ and $\Gal(\Kur/K)$.

Let $\pi \in \calO_K$ be a uniformizer and let $L_\pi$ be the field $K_\pi \Kur$. We define a group homomorphism
$r_\pi : K^\times \to \Gal(L_\pi / K)$ as follows
\begin{enumerate}
    \item $r_\pi (\pi)$ acts as identity on $K_\pi$ and acts as Frobenius automorphism $\sigma$ on
        $\Kur$.
    \item Given a $u \in U_K$, we define $r_\pi (u)$ to act as $[u^{-1}]_f$ on $K_\pi$ and
        as the identity on $\Kur$.
\end{enumerate}

%%%%%%%%%%%%%%%%%%%%%%%%%%%%%%%%%%%%%%%%%%%
% Next, we will prove that the field $L_\pi$ and the homomorphism $r_\pi$ are independent of the
% uniformizer $\pi$, and that $L_\pi$ is a maximal abelian extension of $K$.

% % Moreover, $r_\pi$ coincides with
% % the local reciprocity map in local class field theory, i.e. for all $a \in \calO_K$,
% % $r_\pi(a) = (a, L_\pi \Kur / K ^ \times)$. The references of these can be found in
% % \cite{LubinTate1965} and \cite{cassels1967algebraic}.

% To prove thess results, we first prove two auxiliary lemmas.

% \begin{lemma}
%     Let $E$ be an algebraic extension (finite or infinite) of a local field
%     and let $\alpha \in \widehat{E}$, where $\widehat{E}$ is the completion of $E$.
%     Then, if $\alpha$ is separable algebraic
%     over $E$, then $\alpha$ belongs to $E$.

%     \label{lem: sep_mem}
% \end{lemma}

% \begin{proof}
%     Let $\Fsep{E}$ be the separable closure of $E$ and let $E'$ be the adherence
%     (topological closure) of $E$ in $\Fsep{E}$. Then we can view $\alpha$ as
%     an element in $E'$, as $\alpha \in \widehat{E}$. Hence, it is enough to
%     show that $E' = E$.

%     Let $g \in \Gal(\Fsep{E} / E)$. Since $g$ is continuous and $g$ is identity
%     one $E$, it follows that $g$ is identity on $E'$. Hence, we have that
%     \[\Gal(\Fsep{E}/ E) = \Gal(\Fsep{E}/ E'). \]
%     Then we conclude that $E' = E$.
% \end{proof}

% Let $\sigma \in \Gal(\Kur/ K)$ be the Frobenius automorphism. We claim
% that $\sigma$ can be extended to $\widehat{\Kur}$ by continuity.

% \begin{proof}[Proof of the claim: ]
% Let $v$ denote the normalized valuation on $K$. Since $K$ is complete,
% $v$ extends uniquely to every algebraic extension of $K$, hence in particular
% to $K^{\mathrm{ur}}$. Therefore every $K$-automorphism of $K^{\mathrm{ur}}$
% preserves this valuation. Thus, for all $x,y\in K^{\mathrm{ur}}$,
% \[ v(\sigma(x)-\sigma(y)) = v(\sigma(x-y)) = v(x-y).
% \]
% Equivalently, $\sigma$ is an isometry for the nonarchimedean absolute value
% on $K^{\mathrm{ur}}$.

% Hence $\sigma$ is uniformly continuous. Since $K^{\mathrm{ur}}$ is dense in
% its completion $\widehat{K^{\mathrm{ur}}}$, the map $\sigma$ extends uniquely
% by continuity to a map
% \[
% \widehat{\sigma}:\widehat{K^{\mathrm{ur}}}\longrightarrow
% \widehat{K^{\mathrm{ur}}}.
% \]

% The inverse $\sigma^{-1}$ is also an isometry, so it extends as well, and its
% extension is inverse to $\widehat{\sigma}$. Therefore $\widehat{\sigma}$ is an
% automorphism of $\widehat{K^{\mathrm{ur}}}$ extending $\sigma$.
% For convention, we will denote $\widehat{\sigma}$ as $\sigma$.
% \end{proof}

% Let $\pi$ be a uniformizer of $\calO_K$ and $\omega$ be another uniformizer,
% i.e. $\omega = \pi u$ with $u \in U_K$. Let $f \in \calF_\pi$ and
% $g \in \calF_\omega$.

% \begin{lemma}
%     Let $\sigma \in \Gal(\Kur/ K)$ be the Frobenius automorphism and it can
%     be extended to $\widehat{\Kur}$. Then there is a power series
%     $\phi(X) \in \widehat{\calO_{\Kur}}[\![X]\!]$ with
%     $\phi(X) \equiv \epsilon X \pmod{\deg 2}$ with
%     $\epsilon$ a unit in $\calO_{\Kur}$
%     such that
%     \begin{align}
%         \leftsup{\sigma}{\phi}& = \phi \circ [u]_f, \label{eq: phi_u} \\
%         \phi(F_f(X,Y)) &= F_g(\phi(X), \phi(Y)), \label{eq: phi_hom_fg} \\
%         \phi \circ [a]_f &= [a]_g \circ \phi, \quad \forall a \in \calO_K \label{eq: phi_a} .
%     \end{align}

%     Here, $\leftsup{\sigma}{\phi}$ is the power series obtained by
%     $\sigma$ acting on each coefficient of $\phi$.

%     \label{lem: phi_lemma}
% \end{lemma}

% \begin{proof}
%     We first construct a power series $\varphi(X) \in\widehat{\calO_{\Kur}}[\![X]\!]$
%     satisfying equation \eqref{eq: phi_u} and modify it to satisfy
%     equation \eqref{eq: phi_hom_fg} and equation \eqref{eq: phi_a}.

%     Recall that $\sigma - 1$ is surjective on the additive group
%     of $\widehat{\calO_{\Kur}}$ and it is also surjective
%     on the group of units in $\widehat{\calO_{\Kur}}$. Then there
%     is an element $\epsilon$ which is a unit in $\widehat{\calO_{\Kur}}$
%     such that $\sigma \epsilon = \epsilon u$.

%     Let $\varphi_1 (X) = \epsilon X$, then we have that
%     \[\leftsup{\sigma}{\varphi_1}(X) \equiv \sigma \epsilon X
%         \equiv \epsilon u X \equiv \varphi_1 \circ [u]_f (X) \pmod{\deg 2}.   \]

%     Suppose that
%     \[\leftsup{\sigma}{\varphi_n} (X) \equiv \varphi_n \circ [u]_f (X)
%     \pmod{\deg (n+1)}.   \]

%     Then we claim that $\varphi_{n+1}(X) \coloneqq \varphi_n(X)
%     + a \epsilon^{n +1} X ^ {n +1}$, where $a$ satisfies
%     \[a - \sigma a = \frac{c}{(\sigma \epsilon)^{n+1}}, \]

%     and $c$ is the $n+1$-th coefficient of
%     power series $\leftsup{\sigma}{\varphi_n} (X) - \varphi_n \circ [u]_f (X)$.
%     The $a$ always exists since $\sigma - 1$ is surjective. Now we prove by
%     calculation the $\varphi_{n+1}(X)$ satisfies
%     \begin{equation}
%         \leftsup{\sigma}{\varphi_{n+1}}(X) \equiv \varphi_{n+1} \circ [u]_f (X)
%     \pmod{\deg (n + 2)}.
%     \label{eq: ind_lem_phi}
%     \end{equation}

%     Since $\varphi_{n+1}(X) \coloneqq \varphi_n(X) + a \epsilon ^{n+1} X ^ {n+1}$,
%     it is enough to check that the $n+1$-th coefficient of both
%     side coincide.

%     On the one hand, the $n+1$-th coefficient of $\leftsup{\sigma}{\varphi_{n+1}}(X)$
%     is
%     \[ \text{$n+1$-th coefficient of $\leftsup{\sigma}{\varphi_{n}}(X)$}
%     + \sigma a \sigma (\epsilon^{n+1}).  \]

%     On the other hand, the $n+1$-th coefficient of $\varphi_{n+1} \circ [u]_f (X)$ is
%     \[\text{$n+1$-th coefficient of $\varphi_{n} \circ [u]_f (X)$}
%     + a \epsilon^{n+1} u^{n+1}.  \]

%     Then the equation \eqref{eq: ind_lem_phi} follows from
%     the definition of $a$ and the fact $\sigma \epsilon = \epsilon u$.
%     By induction, equation \ref{eq: ind_lem_phi} holds for all $n$. We obtain
%     a power series $\varphi \coloneqq \lim \varphi_n$ over
%     $\widehat{\calO_{\Kur}}$. Since $\varphi (X) \equiv \epsilon X \pmod{\deg 2}$,
%     there is a substitution inverse of $\varphi$ and we denote it as $\varphi^{-1}$.

%     Consider the power series
%     \[h \coloneqq \leftsup{\sigma}{\varphi} f \varphi^{-1}. \]
%     Then we have that
%     \[h = \leftsup{\sigma}{\varphi} f \varphi^{-1} =
%     \varphi [u]_f f \varphi^{-1} = \varphi [u]_f [\pi]_f \varphi^{-1}
%     = \varphi [\omega]_f \varphi^{-1}.
%     \]
%     The power series $h$ has its coefficients in $\calO_K$, because
%     \[
%     \leftsup{\sigma}{h} = \leftsup{\sigma}{\varphi} \leftsup{\sigma}{[\omega]_f}
%     \leftsup{\sigma}{\varphi^{-1}} =
%     \leftsup{\sigma}{\varphi} [\omega]_f
%     \leftsup{\sigma}{\varphi^{-1}} =
%     \leftsup{\sigma}{\varphi} f [u]_f
%     \leftsup{\sigma}{\varphi^{-1}} = \leftsup{\sigma}{\varphi} f \varphi^{-1} = h.
%     \]
%     where the second-to-last equality use the fact that
%     \[\varphi \circ  [u]_f \leftsup{\sigma}{\varphi}^{-1} =
%     \leftsup{\sigma}{\varphi} \leftsup{\sigma}{\varphi}^{-1} = X, \]
%     and the uniqueness of substitution inverse.

%     Moreover, we have that
%     \[h(X) \equiv \epsilon \omega \epsilon^{-1} X \equiv \omega X \pmod{\deg 2},  \]
%     and
%     \[h(X) \equiv \leftsup{\sigma}{\varphi}(f (\varphi^{-1}(X))) \equiv
%     \leftsup{\sigma}{\varphi}(\varphi^{-1}(X)^q) \equiv
%     \leftsup{\sigma}{\varphi}(\leftsup{\sigma}{\varphi}^{-1}(X^q)) \pmod{\pi}.
%     \]

%     Hence, we have that $h \in \calF_\omega$.
%     Now, we define $\phi$ to be $[1]_{g,h} \varphi$. And we check that this $\phi$
%     satisfies three equations in the lemma. Equation \eqref{eq: phi_u} follows from
%     the equality
%     $\leftsup{\sigma}{[1]_{g,h}} = [1]_{g,h}$.

%     Note that $g = \phi [\omega]_f \phi^{-1}$. It is proved as follow:
%     \[g = [1]_{g,h} h [1]_{g,h}^{-1} =
%     \leftsup{\sigma}{[1]_{g,h}} \leftsup{\sigma}{\varphi} f \varphi^{-1} [1]_{g,h}^{-1}
%     = \leftsup{\sigma}{\phi} f \phi^{-1} =\phi [u]_f f \phi^{-1}
%     = \phi [\omega]_f \phi^{-1}.
%     \]

%     To show the equation \eqref{eq: phi_hom_fg}, we verify that
%     the power series
%     \[F(X,Y) \coloneqq \phi (F_f(\phi^{-1}(X), \phi^{-1}(Y))) =
%     \phi F_f \phi^{-1} \]
%     satisfies the same properties characterizing $F_g(X,Y)$ in Lemma
%     \ref{lemma: constructive_lemma_lubin_Tate}. We need to
%     prove that $F(X,Y) \equiv X + Y \pmod{\deg 2}$ and $g (F(X,Y)) = F(g(X), g(Y))$.
%     The first equality follows by computing coefficient of $X$ and $Y$.
%     The second follows from the following computation
%     \[gF = \phi [\omega]_f \phi^{-1} \phi F_f \phi^{-1}
%     = \phi [\omega]_f F_f \phi^{-1} =
%     \phi F_f [\omega]_f \phi^{-1} = \phi F_f \phi^{-1} (\phi [\omega]_f \phi^{-1})
%     = \phi F_f \phi^{-1} g= F g.
%     \]

%     For equation \eqref{eq: phi_a}, it is enough to verify $\phi [a]_f \phi^{-1}$
%     satisfies the characterizing properties of $[a]_g$ in
%     Lemma \ref{lemma: constructive_lemma_lubin_Tate}. And this follows from
%     \[ \phi [a]_f \phi^{-1} g =\phi [a]_f \phi^{-1} \phi [\omega]_f \phi^{-1}
%     = \phi [a]_f [\omega]_f \phi^{-1} = \phi [a \omega]_f \phi^{-1} =
%     \phi [\omega]_f \phi^{-1} \phi [a]_f\phi^{-1} = g \phi [a]_f\phi^{-1}. \]

%     We conclude that this $\phi$ satisfies all three equation.
% \end{proof}

% \begin{theorem}
%     Let $\omega \in \calO_K$ be another uniformizer; that is,
%     $\omega = u \pi$ for some $u \in U_K$. Then
%     \[ L_\pi = L_\omega. \]
% \end{theorem}

% \begin{proof}
%     Let $f\in \calO_\pi$ and $g \in \calO_\omega$.
%     Then $L_\pi = K (\Lambda_f)$ and $L_\omega = K(\Lambda_g)$.
%     By Lemma \ref{lem: phi_lemma}, there is a power series $\phi$ such that
%     $\phi$ is a isomorphism from formal $\calO_K$-module $\Lambda_f$ to
%     formal $\calO_K$-module $\Lambda_g$.
%     Then we have that $\phi (\Lambda_f) = \Lambda_g$.
%     It follows that $\Lambda_g = \phi(\Lambda_f) \subset
%     \widehat{\Kur (\Lambda_f)} = \widehat{L_\pi}$. By symmetry, we have that
%     $\widehat{L_\pi} = \widehat{L_\omega}$. And the the theorem
%     follows from Lemma \ref{lem: sep_mem}.
% \end{proof}

% \begin{theorem}
%     Let $\omega \in \calO_K$ be another uniformizer; that is,
%     $\omega = u \pi$ for some $u \in U_K$. Then $r_\pi = r_\omega$.
% \end{theorem}

% \begin{proof}
%     Let $f\in \calO_\pi$ and $g \in \calO_\omega$.
%     Then $L_\pi = K (\Lambda_f)$ and $L_\omega = K(\Lambda_g)$.

%     It is enough to show that $r_\pi(\omega) = r_\omega(\omega)$, since $\omega$
%     generates $K^\times$. On the field $\Kur$, the automorphism
%     $r_\pi(\omega)$ and $r_\omega(\omega)$ acts as Frobenius automorphism.
%     It is enough to show that their actions on $L_\omega$ are the same.
%     By definition, $r_\omega(\omega)$ is identity on $L_\omega$.

%     By Lemma \ref{lem: phi_lemma}, there is a power series $\phi$ over
%     $\widehat{\calO_{\Kur}}$ such that $\phi$ is a isomorphism from
%     formal $\calO_K$-module $\Lambda_f$ to
%     formal $\calO_K$-module $\Lambda_g$. Then $\phi(\Lambda_f) = \Lambda_g$.
%     Since $L_\omega = K (\Lambda_g)$, it is enough to show that
%     \[ \leftsup{r_\pi(\omega)}{\phi(\lambda)} = \phi(\lambda), \quad
%     \forall \lambda \in \Lambda_f.  \]

%     Now, we write $r_\pi(\omega) = r_\pi(u) r_\pi(\pi)$. Denote $r_\pi(u)$ as
%     $\tau$ and $r_\pi(\pi)$ as $\sigma$.
%     And by definition of $r_\pi$,
%     $\sigma$ is Frobenius on $\widehat{\Kur}$ and identity on $\lambda$,
%     whereas $\tau$ is identity on $\widehat{\Kur}$ and maps $\lambda$ to
%     $[u^{-1}]_f(\lambda)$. Note that $\phi \in \widehat{\calO_{\Kur}}[\![X]\!]$,

%     \[\leftsup{r_\pi(\omega)}{\phi(\lambda)} =
%     \leftsup{\tau \sigma}{\phi(\lambda)} =
%     \leftsup{\sigma}{\phi}(\tau \lambda) =
%     \phi ([u]_f [u^{-1}]_f (\lambda))
%     = \phi(\lambda).  \]


% \end{proof}
%%%%%%%%%%%%%%%%%%%%%%%%%%%%%

Next, we prove that the field $L_\pi$ and the homomorphism $r_\pi$ are
independent of the choice of uniformizer $\pi$.
Moreover, $L_\pi$ is a maximal abelian extension of $K$.

To prove these results, we first prove two auxiliary lemmas.

\begin{lemma}[{\cite[Chap.~VI 3.3 Lemma 2]{cassels1967algebraic}}]
    Let $E$ be an algebraic extension, finite or infinite, of a local field, and
    let $\alpha \in \widehat{E}$, where $\widehat{E}$ denotes the completion of
    $E$. If $\alpha$ is separable algebraic over $E$, then $\alpha \in E$.
    \label{lem: sep_mem}
\end{lemma}

\begin{proof}
    Let $\Fsep{E}$ be the separable closure of $E$, and let $E'$ be the
    adherence, that is, the topological closure of $E$ in $\Fsep{E}$.
    Since $\alpha$ is separable algebraic over $E$, after choosing an
    $E$-embedding of $E(\alpha)$ into $\Fsep{E}$, we may view $\alpha$ as
    an element of $E'$. Hence it is enough to show that $E'=E$.

    Let $g \in \Gal(\Fsep{E}/E)$. Since $g$ is continuous and is the identity
    on $E$, it is also the identity on the closure $E'$. Therefore
    $\Gal(\Fsep{E}/E)=\Gal(\Fsep{E}/E')$. By the fundamental theorem of
    Galois theory, applied to the Galois extension $\Fsep{E}/E$, the fixed
    fields of these two groups are equal. Thus $E'=E$.
\end{proof}

Let $\sigma \in \Gal(\Kur/K)$ be the Frobenius automorphism. We claim that
$\sigma$ extends uniquely to $\widehat{\Kur}$ by continuity.

\begin{proof}[Proof of the claim]
    Let $v$ denote the normalized valuation on $K$. Since $K$ is complete,
    $v$ extends uniquely to every algebraic extension of $K$, and hence in
    particular to $\Kur$. Therefore every $K$-automorphism of $\Kur$ preserves
    this valuation. Thus, for all $x,y \in \Kur$,
    \[
        v(\sigma(x)-\sigma(y)) = v(\sigma(x-y)) = v(x-y).
    \]
    Equivalently, $\sigma$ is an isometry for the corresponding
    nonarchimedean absolute value.

    Hence $\sigma$ is uniformly continuous. Since $\Kur$ is dense in
    $\widehat{\Kur}$, the map $\sigma$ extends uniquely to a continuous map
    $\widehat{\sigma}:\widehat{\Kur}\to \widehat{\Kur}$. The inverse
    $\sigma^{-1}$ is also an isometry, so it extends as well, and its extension
    is inverse to $\widehat{\sigma}$. Thus $\widehat{\sigma}$ is an
    automorphism of $\widehat{\Kur}$ extending $\sigma$. By abuse of notation,
    we shall again denote this extension by $\sigma$.
\end{proof}

Let $\pi$ be a uniformizer of $\calO_K$, and let $\omega$ be another
uniformizer. Write $\omega=u\pi$ with $u \in U_K$. Let
$f \in \calF_\pi$ and $g \in \calF_\omega$.

\begin{lemma}[{\cite[Lemma 2]{LubinTate1965}}]
    There exists a power series
    $\phi(X)$ over $\widehat{\calO_{\Kur}}$ such that
    $\phi(X) \equiv \epsilon X \pmod{\deg 2}$ for some unit
    $\epsilon \in \widehat{\calO_{\Kur}}^\times$, and such that
    \begin{align}
        \leftsup{\sigma}{\phi} &= \phi \circ [u]_f,
            \label{eq: phi_u} \\
        \phi(F_f(X,Y)) &= F_g(\phi(X),\phi(Y)),
            \label{eq: phi_hom_fg} \\
        \phi \circ [a]_f &= [a]_g \circ \phi,
            \qquad \forall a \in \calO_K .
            \label{eq: phi_a}
    \end{align}
    Here $\leftsup{\sigma}{\phi}$ denotes the power series obtained by applying
    $\sigma$ to each coefficient of $\phi$.
    \label{lem: phi_lemma}
\end{lemma}

\begin{proof}
    We first construct a power series
    $\varphi(X) \in \widehat{\calO_{\Kur}}[\![X]\!]$ satisfying
    \eqref{eq: phi_u}. We will then modify it so that it also satisfies
    \eqref{eq: phi_hom_fg} and~\eqref{eq: phi_a}.

    Recall that $\sigma-1$ is surjective on the additive group of
    $\widehat{\calO_{\Kur}}$, and also on the group
    $\widehat{\calO_{\Kur}}^\times$ of units. Hence there exists a unit
    $\epsilon \in \widehat{\calO_{\Kur}}^\times$ such that
    $\leftsup{\sigma}{\epsilon}=\epsilon u$.

    Set $\varphi_1(X)=\epsilon X$. Then
    $\leftsup{\sigma}{\varphi_1}(X) \equiv
    \varphi_1([u]_f(X)) \pmod{\deg 2}$.

    Suppose inductively that
    \[
        \leftsup{\sigma}{\varphi_n}(X)
        \equiv
        \varphi_n([u]_f(X))
        \pmod{\deg(n+1)}.
    \]
    Let $c$ be the coefficient of $X^{n+1}$ in
    $\leftsup{\sigma}{\varphi_n}(X)-\varphi_n([u]_f(X))$. Choose
    $a \in \widehat{\calO_{\Kur}}$ such that
    \[
        a-\leftsup{\sigma}{a}
        =
        \frac{c}{(\leftsup{\sigma}{\epsilon})^{n+1}},
    \]
    which is possible by the surjectivity of $\sigma-1$. Define
    \[
        \varphi_{n+1}(X)
        =
        \varphi_n(X)+a\epsilon^{n+1}X^{n+1}.
    \]
    Since $\leftsup{\sigma}{\epsilon}=\epsilon u$, the coefficient of
    $X^{n+1}$ in
    $\leftsup{\sigma}{\varphi_{n+1}}(X)-\varphi_{n+1}([u]_f(X))$
    is
    \[
        c+\leftsup{\sigma}{a}(\leftsup{\sigma}{\epsilon})^{n+1}
        -a\epsilon^{n+1}u^{n+1}
        =
        c-\bigl(a-\leftsup{\sigma}{a}\bigr)
        (\leftsup{\sigma}{\epsilon})^{n+1}
        =
        0.
    \]
    Therefore
    \[
        \leftsup{\sigma}{\varphi_{n+1}}(X)
        \equiv
        \varphi_{n+1}([u]_f(X))
        \pmod{\deg(n+2)}.
    \]
    Passing to the coefficient-wise limit, we obtain a power series
    $\varphi=\lim_n \varphi_n$ satisfying
    $\leftsup{\sigma}{\varphi}=\varphi\circ [u]_f$. Since
    $\varphi(X)\equiv \epsilon X \pmod{\deg 2}$ with $\epsilon$ a unit,
    $\varphi$ has a substitution inverse, denoted by $\varphi^{-1}$.

    Now define
    \[
        h
        =
        \leftsup{\sigma}{\varphi}\circ f\circ \varphi^{-1}.
    \]
    Since $\leftsup{\sigma}{\varphi}=\varphi\circ [u]_f$ and
    $f=[\pi]_f$, we have
    \[
        h
        =
        \varphi\circ [u]_f\circ [\pi]_f\circ \varphi^{-1}
        =
        \varphi\circ [\omega]_f\circ \varphi^{-1}.
    \]
    Applying $\sigma$ to this equality gives
    \[
        \leftsup{\sigma}{h}
        =
        \leftsup{\sigma}{\varphi}\circ [\omega]_f
        \circ (\leftsup{\sigma}{\varphi})^{-1}
        =
        \varphi\circ [u]_f\circ [\omega]_f\circ [u^{-1}]_f
        \circ \varphi^{-1}
        =
        h.
    \]
    Hence the coefficients of $h$ are fixed by $\sigma$, and therefore lie in
    $\calO_K$.

    Moreover,
    $h(X)\equiv \omega X \pmod{\deg 2}$. Also, using
    $f(X)\equiv X^q \pmod{\pi}$ and the fact that $\sigma$ lifts the
    $q$-power Frobenius on the residue field, we get
    \[
        h(X)
        =
        \leftsup{\sigma}{\varphi}(f(\varphi^{-1}(X)))
        \equiv
        \leftsup{\sigma}{\varphi}\bigl((\varphi^{-1}(X))^q\bigr)
        \equiv
        \leftsup{\sigma}{\varphi}
        \bigl((\leftsup{\sigma}{\varphi})^{-1}(X^q)\bigr)
        =
        X^q
        \pmod{\pi}.
    \]
    Since $\pi$ and $\omega$ generate the same maximal ideal, this shows that
    $h \in \calF_\omega$.

    We now replace $\varphi$ by
    \[
        \phi=[1]_{g,h}\circ \varphi.
    \]
    Since $[1]_{g,h}$ has coefficients in $\calO_K$, it is fixed by $\sigma$.
    Thus
    \[
        \leftsup{\sigma}{\phi}
        =
        [1]_{g,h}\circ \leftsup{\sigma}{\varphi}
        =
        [1]_{g,h}\circ \varphi\circ [u]_f
        =
        \phi\circ [u]_f,
    \]
    so~\eqref{eq: phi_u} still holds.

    Next, since $[1]_{g,h}$ is an isomorphism from $F_h$ to $F_g$, it satisfies
    $g\circ [1]_{g,h}=[1]_{g,h}\circ h$. Hence
    \[
        \leftsup{\sigma}{\phi}\circ f\circ \phi^{-1}
        =
        [1]_{g,h}\circ
        \leftsup{\sigma}{\varphi}\circ f\circ \varphi^{-1}
        \circ [1]_{g,h}^{-1}
        =
        [1]_{g,h}\circ h\circ [1]_{g,h}^{-1}
        =
        g.
    \]
    Combining this with~\eqref{eq: phi_u}, we also have
    $g=\phi\circ [\omega]_f\circ \phi^{-1}$.

    To prove~\eqref{eq: phi_hom_fg}, define
    \[
        F(X,Y)
        =
        \phi\bigl(F_f(\phi^{-1}(X),\phi^{-1}(Y))\bigr).
    \]
    Then $F(X,Y)\equiv X+Y \pmod{\deg 2}$.
    Here and below, for any one-variable power series $h$, we abuse notation by
writing $hF$ for $h\circ F$ and $Fh$ for $F(h(X),h(Y))$.
    Moreover, using
    $g=\phi\circ [\omega]_f\circ \phi^{-1}$ and the fact that $[\omega]_f$
    is an endomorphism of $F_f$, we have
    % \[
    %     g(F(X,Y))
    %     =
    %     F(g(X),g(Y)).
    % \]
    \[gF = \phi [\omega]_f \phi^{-1} \phi F_f \phi^{-1}
    = \phi [\omega]_f F_f \phi^{-1} =
    \phi F_f [\omega]_f \phi^{-1} = \phi F_f \phi^{-1} (\phi [\omega]_f \phi^{-1})
    = \phi F_f \phi^{-1} g= F g.
    \]
    By the uniqueness statement in Lemma~\ref{lemma: constructive_lemma_lubin_Tate},
    it follows that $F=F_g$. This proves~\eqref{eq: phi_hom_fg}.

    Finally, fix $a\in \calO_K$ and set
    $A_a=\phi\circ [a]_f\circ \phi^{-1}$. Then
    $A_a(X)\equiv aX \pmod{\deg 2}$. Also,
    \[
        A_a\circ g
        =
        \phi\circ [a]_f\circ [\omega]_f\circ \phi^{-1}
        =
        \phi\circ [\omega]_f\circ [a]_f\circ \phi^{-1}
        =
        g\circ A_a.
    \]
    Again by the uniqueness statement in
    Lemma~\ref{lemma: constructive_lemma_lubin_Tate}, we get
    $A_a=[a]_g$. Hence~\eqref{eq: phi_a} holds.
\end{proof}

\begin{theorem}[{\cite[Theorem 3]{LubinTate1965}}]
    The field $L_\pi$ is independent of the choice of uniformizer $\pi$.
\end{theorem}

\begin{proof}
    In this proof, write
    \[
        L_\pi=\Kur(\Lambda_f)
        \qquad\text{and}\qquad
        L_\omega=\Kur(\Lambda_g),
    \]
    where $f\in \calF_\pi$ and $g\in \calF_\omega$.

    By Lemma~\ref{lem: phi_lemma}, there exists a power series $\phi$ giving
    an isomorphism from the formal $\calO_K$-module associated with $f$ to the
    one associated with $g$. In particular, $\phi(\Lambda_f)=\Lambda_g$.

    Since the coefficients of $\phi$ lie in $\widehat{\calO_{\Kur}}$, for every
    $\lambda\in \Lambda_f$ we have
    $\phi(\lambda)\in \widehat{\Kur(\Lambda_f)}=\widehat{L_\pi}$. Thus
    $\Lambda_g\subset \widehat{L_\pi}$. Each element of $\Lambda_g$ is
    separable algebraic over $L_\pi$, so Lemma~\ref{lem: sep_mem} implies that
    $\Lambda_g\subset L_\pi$. Therefore $L_\omega\subset L_\pi$.

    By symmetry, $L_\pi\subset L_\omega$. Hence $L_\pi=L_\omega$.
\end{proof}

\begin{theorem}[{\cite[Theorem 3]{LubinTate1965}}]
    The homomorphism $r_\pi$ is independent of the choice of uniformizer
    $\pi$.
\end{theorem}

\begin{proof}
    It is enough to prove the following assertion: if $\omega=u\pi$ is another
    uniformizer, then
    \[
        r_\pi(\omega)=r_\omega(\omega).
    \]
    Indeed, applying this assertion with $\omega$ replaced by an arbitrary
    uniformizer shows that all the maps $r_\pi$ have the same value on every
    uniformizer. Since the set of uniformizers generates $K^\times$, the
    homomorphisms $r_\pi$ are equal.

    Fix $f\in \calF_\pi$ and $g\in \calF_\omega$. By the previous theorem,
    we may identify
    \[
        L_\pi=L_\omega=\Kur(\Lambda_g).
    \]
    On $\Kur$, both $r_\pi(\omega)$ and $r_\omega(\omega)$ induce the
    Frobenius automorphism. Hence it remains to compare their actions on
    $\Lambda_g$. By definition, $r_\omega(\omega)$ is the identity on
    $\Lambda_g$.

    By Lemma~\ref{lem: phi_lemma}, there exists a power series $\phi$ such
    that $\phi(\Lambda_f)=\Lambda_g$ and
    $\leftsup{\sigma}{\phi}=\phi\circ [u]_f$. Therefore it is enough to prove
    that, for every $\lambda\in \Lambda_f$,
    \[
        \leftsup{r_\pi(\omega)}{\phi(\lambda)}
        =
        \phi(\lambda).
    \]

    Write $r_\pi(\omega)=r_\pi(u)r_\pi(\pi)$, and set
    $\tau=r_\pi(u)$ and $\sigma=r_\pi(\pi)$. By definition, $\sigma$ acts as
    Frobenius on $\widehat{\Kur}$ and fixes $\Lambda_f$, while $\tau$ fixes
    $\widehat{\Kur}$ and sends $\lambda\in\Lambda_f$ to
    $[u^{-1}]_f(\lambda)$. Since the coefficients of $\phi$ lie in
    $\widehat{\calO_{\Kur}}$, we get
    \[
        \leftsup{r_\pi(\omega)}{\phi(\lambda)}
        =
        \leftsup{\tau\sigma}{\phi(\lambda)}
        =
        \leftsup{\sigma}{\phi}(\leftsup{\tau}{\lambda})
        =
        \phi([u]_f([u^{-1}]_f(\lambda)))
        =
        \phi(\lambda).
    \]
    Thus $r_\pi(\omega)=r_\omega(\omega)$, and the theorem follows.
\end{proof}


\begin{corollary}[{\cite[Corollary]{LubinTate1965}}]
    Let $\pi$ be a uniformizer of $\calO_K$. Then $L_\pi\Kur$ is the maximal
    abelian extension of $K$, and $r_\pi$ is the local reciprocity homomorphism
    for $L_\pi\Kur/K$. More precisely, for every $a\in K^\times$,
    $r_\pi(a)=(a,L_\pi\Kur/K)$, with respect to our normalization of the
    reciprocity map.
\end{corollary}

\begin{proof}
    By the previous theorem, $r_\pi$ is independent of the choice of
    uniformizer. Thus we may denote this common homomorphism simply by
    $r:K^\times\to \Gal(L_\pi\Kur/K)$.

    We first explain why the maximality of $L_\pi\Kur$ follows once we know
    that $r$ is the local reciprocity homomorphism. Since $L_\pi\Kur$ contains
    $\Kur$, if $r(a)$ is trivial on $L_\pi\Kur$, then its restriction to $\Kur$
    is trivial. But this restriction is $\sigma^{v_K(a)}$, where $\sigma$ is
    the Frobenius automorphism. Hence $v_K(a)=0$, so $a\in U_K$.

    Now let $u\in U_K$. If $r(u)$ is trivial on $L_{\pi,m}$, then, by the
    definition of $r_\pi$ on units, the endomorphism $[u^{-1}]_f$ acts
    trivially on $\Lambda_{f,m}$. Equivalently, $u\equiv 1 \pmod{\pi^m}$.
    If $r(u)$ is trivial on all of $L_\pi=\bigcup_{m\geq 1}L_{\pi,m}$, then
    this congruence holds for every $m$, and therefore $u=1$. Hence the
    restriction of $r$ to $U_K$ is injective, and consequently $r$ itself is
    injective.

    Therefore, if $r$ is the restriction of the local reciprocity map to
    $L_\pi\Kur$, then the corresponding kernel is trivial. By the Galois
    correspondence in local class field theory, this means that $L_\pi\Kur$
    corresponds to the trivial subgroup of $K^\times$. Hence $L_\pi\Kur$ is
    the maximal abelian extension of $K$.

    It remains to show that $r$ agrees with the local reciprocity map. Let
    $s:K^\times\to \Gal(L_\pi\Kur/K)$ be the local reciprocity homomorphism,
    so that $s(a)=(a,L_\pi\Kur/K)$.

    Let $\pi$ be any uniformizer of $\calO_K$. By Theorem~\ref{cor: norm_Lpi},
    the element $\pi$ is a norm from $L_{\pi,m}/K$ for every $m$. Hence, by the
    defining property of the local reciprocity map, $s(\pi)$ acts trivially on
    $L_\pi=\bigcup_{m\geq 1}L_{\pi,m}$. On the other hand, its restriction to
    $\Kur$ is the Frobenius automorphism $\sigma$.

    By definition, $r_\pi(\pi)$ has exactly the same two properties: it is the
    identity on $L_\pi$, and it acts as $\sigma$ on $\Kur$. Therefore
    $s(\pi)=r_\pi(\pi)$. Since $r_\pi=r$ by the independence of the uniformizer,
    we get $s(\pi)=r(\pi)$ for every uniformizer $\pi$.

    Finally, the uniformizers generate $K^\times$. Indeed, after fixing one
    uniformizer $\pi$, every element of $K^\times$ can be written as
    $u\pi^n$ with $u\in U_K$. Thus $s$
    and $r$ agree on a set of generators of $K^\times$, and hence $s=r$.
    This proves that $r_\pi$ is the local reciprocity homomorphism.
\end{proof}



\section{Connection with functional equation lemma}

In this section, we discuss the connection between the formal group laws obtained from the
functional equation lemma and Lubin--Tate formal group laws. Throughout this section,
we fix a nonarchimedean local field $K$. Let $\calO_K$ be its ring of integers,
$\pi$ a uniformizer, and $q$ the cardinality of its residue field.

\begin{proposition}
    The Lubin--Tate formal group laws $F_f$ associated to $f\in \calF_\pi$ are in one-to-one
    correspondence with formal group laws obtained by the functional equation with
    parameters $I = (\pi), \sigma = \id, s_1 = \pi^{-1}$ and $\forall i \ge 2, s_i = 0$.
    The functional equation in the sense of these parameters is
    \[ f_g(X) = g(X) + \pi^ {-1}f_g(X ^ q). \]
    This correspondence is given by
    \[g(X) \mapsto f_g^{-1}(\pi f_g(X)) \in \calF_\pi. \]
\end{proposition}

\begin{proof}
    The integrality of the coefficients in $f_g^{-1}(\pi f_g(X))$ follows from the
    functional equation lemma. Multiplying both sides of the functional equation by $\pi$,
    we have that
    \[ \pi f_g(X) = \pi g(X) + \pi^{-1} \pi f_g(X). \]
    Then $\pi f_g(X) = f_{\pi g}(X)$. According to the second part of the functional
    equation lemma, we have that
    \[f_g^{-1} (\pi f_g(X)) = f_g^{-1} (f_{\pi g}(X)) \in \calO_K[\![X]\!]. \]
    The formal group law obtained by the functional equation lemma is
    \[ F(X,Y) = f_g^{-1} (f_g(X) + f_g(Y)) . \]
    By the existence and uniqueness in Theorem~\ref{thm: Lubin_Tate}, it is enough
    to check that $f_g^{-1}(\pi f_g(X))$ is an endomorphism for $F(X,Y)$. This can be done as follows:
    \begin{align*}
        f_g^{-1} (\pi f_g(F(X,Y))) &= f_g^{-1} (\pi f_g(f_g^{-1} (f_g(X) + f_g(Y)))) \\
            & = f_g^{-1} (\pi f_g(X) + \pi f_g(Y)) \\
            & = F(f_g^{-1}(\pi f_g(X)), f_g^{-1}(\pi f_g(Y))).
    \end{align*}
\end{proof}